% !TEX root = ../main.tex
% !TEX program = lualatex

\documentclass[../main.tex]{subfiles}
\externaldocument{../main}

\begin{document}

\section{Integrated Case Study}

\subsection{Purpose and scope}

The Integrated Case Study (ICS) is an EDO School application exercise, not an
industrial-control-system topic.  Coursebook 3.8.1-4 uses the case study to
make students apply program-office organization, requirements, funding,
technical authority, systems engineering, and stakeholder communication to an
integrated acquisition problem
\autocite{edo-3-8-1-ics-introduction-2025}.

\subsection{Expected products}

\begin{longtblr}[
  caption = {Integrated Case Study products and board value.},
  label = {tab:integrated-case-study-products},
]{
  width = \linewidth,
  colspec = {X[0.9,l] X[1.65,l]},
  rowhead = 1,
  hlines,
  vlines,
}
\textbf{Product} & \textbf{Purpose} \\
Team charter & Defines team purpose, roles, norms, decision process, and delivery expectations. \\
Stakeholder map & Identifies who owns requirements, funding, technical authority, fleet need, and execution risk. \\
Reporting chain & Shows how the issue moves through SYSCOM, PEO, program office, resource sponsor, and fleet channels. \\
Requirements trace & Connects the proposed action to capability need, acquisition phase, and technical baseline. \\
Funding trace & Identifies the appropriation, fiscal year, color-of-money risk, and execution constraint. \\
Presentation package & Forces the team to defend the recommendation with evidence, assumptions, risks, and alternatives. \\
\end{longtblr}

\subsection{Team and IPT behavior}

The case study is graded as both technical reasoning and team performance.  A
strong team makes assumptions explicit, separates facts from judgment, assigns
owners, controls scope, and resolves disagreements through evidence.  Common
leadership barriers include unclear decision authority, unspoken assumptions,
failure to include the right stakeholders, and treating a symptom as the root
problem
\autocite{edo-3-8-1-ics-introduction-2025}.

\subsection{Program-office application checklist}

Use this checklist to translate the case-study ground rules and stakeholder
requirements into a defensible program-office recommendation
\autocite{edo-3-8-1-ics-introduction-2025}.

\begin{itemize}
  \item State the operational problem and requirement source.
  \item Identify the owning program office, resource sponsor, technical authority, and fleet stakeholder.
  \item Define the acquisition phase, decision point, and required artifact.
  \item Trace funding to appropriation, fiscal year, and execution constraint.
  \item Identify technical, schedule, cost, test, logistics, cybersecurity, and contracting risks.
  \item Present alternatives, recommendation, assumptions, and residual risk.
\end{itemize}

\paragraph{Board cue.}
If a board question sounds like an integrated scenario, answer as a
program-office problem: requirement, authority, funding, technical baseline,
stakeholders, risks, alternatives, and recommendation.

\end{document}
